\documentclass[12pt]{article}
%\usepackage{Sweave}
\usepackage[margin=1in]{geometry}
%\usepackage[parfill]{parskip}
%\usepackage[hidelinks]{hyperref}
\usepackage[colorlinks = true,
            linkcolor = black,
            urlcolor  = blue,
            citecolor = black,
            anchorcolor = black]{hyperref}
\usepackage{rotating}
%\usepackage[backend=biber, style=apa, natbib=true]{biblatex}
\usepackage[citestyle=authoryear, style = apa, natbib=true, backend=bibtex, maxcitenames=2]{biblatex}
\addbibresource{references.bib}
\newtheorem{hyp}{H}
\newtheorem{hypo}{Hypothesis}
\usepackage[compact]{titlesec}
%\usepackage{titlesec}
\titleformat*{\section}{\bfseries\fontsize{14}{18}\selectfont}
\titleformat*{\subsection}{\bfseries\fontsize{12}{16}\selectfont}
%\usepackage{natbib}
\usepackage{csquotes}
\usepackage{xfrac}
\usepackage{mathtools}
\usepackage{booktabs}
\usepackage{lscape}
\usepackage{amsmath}
\usepackage[export]{adjustbox}
\usepackage{graphics}
\usepackage{array}
\usepackage[font=small,skip=1pt]{caption}
\usepackage{subcaption} 
\usepackage{ragged2e}
\usepackage{float}
\usepackage[section]{placeins}
\usepackage{hyperref}
\usepackage{multicol}
\usepackage{setspace}
\usepackage[hang,flushmargin]{footmisc}
\usepackage{afterpage}
\usepackage{soul}
\usepackage{xr}
\usepackage{booktabs}
\usepackage{multirow}
\usepackage{siunitx}
\usepackage{caption}
\usepackage{threeparttable}
\usepackage{makecell} 
\usepackage{tabularx}
\newcommand{\tabitem}{\textbullet~~}
\allowdisplaybreaks

% Set caption to be left-aligned
\captionsetup{justification=raggedright, singlelinecheck=false}

% Configure siunitx
\sisetup{
  group-separator = {,},
  group-minimum-digits = 4 % Start grouping from 4 digits
}

\doublespacing

\title{%
  The Signaling Gap: \\
  \large State-Owned Media and the Strategic Management of Sentiment during the Russian Invasion of Ukraine}
\author{Donald Moratz and Sarah Gerstein}
\date{\today}

%\title{The Signaling Gap: State-Owned Media and the Strategic Management of Sentiment during the Russian Invasion of Ukraine}

\begin{document}

\maketitle


\abstract{Alliance commitments are designed to facilitate cooperation, but they can also prevent states from expressing disagreement through formal diplomatic channels, even when they genuinely dissent. We argue that state-owned media provides a mechanism for constrained states to signal genuine preferences to foreign governments, occupying a middle ground between formal channels that are too visible and private diplomacy that lacks credibility. Analyzing state media coverage before and after the 2022 invasion of Ukraine, we find that the invasion produced a larger negative shift in sentiment toward Russia in states with Russian security agreements than in neutral states. This occurred despite these same states refusing to formally condemn the invasion. Our analysis bridges a gap between international relations scholarship, which has largely treated state media as a domestic tool, and Western intelligence services, which have monitored it as a window into regime preferences since the Cold War.}


%%%%%%%%%%%%%%%%%%%%%%%%%%%%%%%% Second Rewrite %%%%%%%%%%%%%%%%%%%%%%%%%%%%%%%%%%%


%%%%%%%%%%%%%%%%%%%%%%%%%%%%%%%% Rewrite %%%%%%%%%%%%%%%%%%%%%%%%%%%%%%%%%%%

\section*{Literature Review}


The idea that state and independent media might react differently to a violation of an international border is not, on its face, surprising. The two media sources operate under different mandates, face different incentives, and are accountable to different principals. What is less obvious is whether these differences are strategically exploited by states during international crises, and if so, how. Answering this requires drawing on three bodies of literature: alliance theory, which establishes the structural constraints states in our sample face; research on state media as a political instrument, which establishes both the incentives and the demonstrated effectiveness of state media manipulation; and the literature on media and diplomatic signaling, which provides the theoretical foundation for our central argument. We develop our hypotheses formally in the following section; here we explore the literature which drives them.

The states in our sample did not face the Russian invasion of Ukraine from a position of strategic freedom. Post-Soviet states occupy a difficult position due to the structure of the international system. They find themselves caught between competing security and economic pressures from Russia and from Western institutions. Alliance theory helps characterize this position. States generally face a choice between balancing against a threatening power or bandwagoning with the stronger side \citep{waltz_theory_2010, walt_origins_2007}. Walt argues that the decision turns not just on raw capabilities, but on the perceived offensive intentions of potential partners where states are more likely to balance against those they perceive as threatening \citep{walt_origins_2007}. Some post-Soviet states have responded to Russian power by tethering themselves to it, using formal security agreements to mitigate the threat of a powerful neighbor \citep{weitsman_dangerous_2004}. The consequence of tethering, however, is constraint. States that have bound themselves to a regional pole through agreements like the CSTO find their capacity for independent signaling substantially limited, particularly when the patron initiates an action that strains the relationship. This structural constraint is the starting point for our argument; it explains why formal diplomatic channels may be unavailable to states that nonetheless have preferences they wish to communicate.

Once states have decided on alliance partners, signaling commitment to those partners, and, critically, signaling dissatisfaction without triggering retaliation, is a fundamental challenge. Signaling can involve words or deeds \citep{fuhrmann_signaling_2014}, and the most visible public signals are those transmitted through media. The news media sets agendas, primes audiences, and transmits elite cues \citep{iyengar_news_2010, chong_framing_2007}. Public statements broadcast through media can function as hand-tying devices, making commitments more credible by making retraction costly \citep{fearon_signaling_1997}. This creates an incentive for states to use media strategically, but as we argue below, their capacity to do so, and the specific channels available, varies systematically with regime type and alliance structure.

State-owned and private media operate under fundamentally different paradigms. Public choice theory suggests that state-owned media outlets distort and manipulate information to prop up incumbent regimes, while privately owned outlets offer comparatively less biased coverage \citep{djankov_who_2003}. State-owned media may also bear the imprint of prior authoritarian regimes \citep{moehler_whose_2011}, though the relationship is not simple: government affiliation and stable funding can in some contexts produce more consistent, if more constrained, coverage. What matters for our purposes is that state-owned and independent media face fundamentally different incentives, and that this difference is likely to be most visible in moments of crisis when the costs of editorial independence are highest and the value of strategic framing is greatest.

The strategic use of state-owned media has become increasingly sophisticated. Recent work has highlighted the extent to which savvy autocracies maintain power through the strategic manipulation of information---management of state-owned media has become particularly consequential during crises \citep{guriev_spin_2022}. At the same time, there are important differences both across and within autocracies. The underlying resource wealth of a state has been shown to have a profound impact on levels of media freedom \citep{egorov_why_2009}, which points to a need to control for the underlying media environment within a state when comparing coverage across our sample. More broadly, \citet{weeks_autocratic_2008} argues that different autocracies have systematic differences in their behavior in the international sphere. We argue that a result of these systematic differences is variation in their capacity for credible formal signaling. Those states with less capability for formal signaling may have an incentive to supplement their formal signals through informal channels that can convey preferences that formal diplomatic constraints make difficult to express directly. We argue that state-owned media is a plausible channel for this communication for two reasons. First, there is historical reliance by many countries in our sample on state-owned media for international signaling \citep{cohen_theatre_1987}. Second, the explicit rise in efforts by western countries to systematically study state media \citep{manor_digitalization_2019, usic_ic_2024} gives senders reason to believe that the signal will be received and correctly interpreted.

The theoretical basis for why informal channels are particularly valuable for autocracies is grounded in formal signaling theory. Credible diplomatic signals require that the signaling state incur costs that make bluffing irrational---either by tying hands through audience costs or by sinking material costs \citep{fearon_signaling_1997}. Democratic governments can generate audience costs through open political competition, where visible domestic dissent or support provides external observers with reliable information about a government's true resolve \citep{schultz_democracy_2001}. Autocracies and weak democracies, which suppress dissent and coerce apparent consensus, face structural credibility deficits in formal channels: their official diplomatic statements are less informative precisely because external observers know they carry little or no domestic political cost. State media, in this context, becomes a more credible---albeit still deniable---channel because foreign governments monitor it as a less-controlled reflection of elite preferences \citep{bjola_revisiting_2018}. The digitalization of public diplomacy has formalized this dynamic. State media is now systematically tracked by foreign intelligence and diplomatic services as informative about the actual policy preferences of the sending government \citep{manor_digitalization_2019, usic_ic_2024}.

There is strong empirical evidence that state-owned media is both strategically deployed and consequential in our region of interest. \citet{gehlbach_government_2014} provides a workhorse theoretical model for strategic state-owned media. They argue that state-owned media is a strategic tool used by governments to advance a specific agenda, and as a result, bias in state-owned media is generally higher than in independent media. While they argue that, with sufficient incentive to advance their agenda, the government is more likely to engage in capture of independent media, either directly or indirectly, in a crisis like the invasion of Ukraine, either strategy is unlikely to be immediately actionable in the short term. Building upon this logic, we argue that we should expect to see a strong divergence between state-owned media coverage and independent media coverage around the invasion as state-owned media responds strategically and before governments can begin to influence independent media coverage. 

The effectiveness of state media in shaping opinion in this kind of context is empirically documented: \citet{enikolopov_media_2011} show that exposure to independent media dramatically reduces vote-share for the ruling party in the 1999 Russian election, highlighting the extent to which state-owned media framing matters in its absence. They find that media plays a particularly powerful role in persuasion in countries with weak democratic institutions, which characterizes most of the states in our sample. At the regional level, \citet{peisakhin_electoral_2018} use variation in exposure to Russian state television to examine support for pro-Russian candidates in Ukraine's 2014 elections, finding that exposure increases support primarily among citizens already inclined toward those candidates. This provides direct evidence that state-owned media functions as a geopolitical instrument in the post-Soviet space. \citet{oates_television_2006} reinforces this picture at the systemic level, arguing that state-owned media in the region serves most directly as a tool for elites to engage in agenda setting, managing political discourse and short-circuiting democratic contestation--- a capacity that becomes especially important when regimes need to manage public reaction to an external shock like the Ukrainian invasion \citep{oates_russian_2016}.

The relationship between media and foreign policy has been examined through several lenses, though prior work has largely treated media as a constraint on government rather than as an instrument of it. The foundational text, \citet{robinson_cnn_2002}, explored the idea of a ``CNN effect'' whereby media coverage of crises forces government intervention, and ultimately found mixed evidence: media influence is only likely when policy is uncertain and coverage is both critical of government and empathetic toward those suffering. Once policy is settled, media largely follows elite cues. We reframe this argument: rather than simply following elite cues or driving foreign policy, states with access to state-owned media strategically weaponize media as foreign policy. \citet{entman_projections_2004} introduced a cascade model in which frames originate with the executive and filter down through elites, media, and finally to the public, with friction at each level within a democracy. We build on his framework by extending it to autocracies and weak democracies, where there is less friction at each level with respect to state-owned media---driven by editorial processes that explicitly reinforce official state positions and by journalistic incentives that impose potentially harsh costs for deviation from them.

\citet{cohen_theatre_1987} emphasizes the role of diplomatic signaling through the media. He argues that these media signals are ambiguous by design, allowing deniability, but are well understood by the intended recipients---a dynamic that has only intensified as foreign governments now systematically monitor state media output \citep{manor_digitalization_2019, bjola_revisiting_2018, usic_ic_2024}. Indeed, during the Cold War, western analysts meticulously studied the state-owned media of the USSR and its satellites to understand Soviet foreign policy. Soviet officials strategically used this fact to signal to the West in ways that they could not do through formal channels \citep{cohen_theatre_1987}. This work serves as a precursor to \citet{miskimmon_strategic_2013}, which introduces the concept of strategic narratives---frameworks that states leverage through media as a core instrument of foreign policy to shape how they project their narrative to the world. Their framework establishes how actors form and project narratives at three levels: system, identity, and issue.

We extend the framework proposed by \citet{miskimmon_strategic_2013} by considering the constrained nature of international signaling due to a combination of trade and security integrations. We posit that strategic narrative projection operates at both formal levels (UN votes, official diplomatic statements) and more informal levels (state-owned media). We argue that the invasion of Ukraine triggers crises at all three levels of the framework in \citet{miskimmon_strategic_2013}. While the invasion serves foremost as an issue narrative contest, the implications of a major land war in Europe---an international order shaking crisis---trigger both system and identity pressures for post-Soviet states. The invasion is a direct challenge to the established post-Cold War system, challenging norms around the use of force and wars of territorial expansion. At the identity level, it challenged state positions: their values, their alliances, and their place in the world. For post-Soviet states caught between Russian security agreements and Western integration, this introduced enormous identity pressures. That this crisis provoked challenges at all three levels introduced a tension that states needed to manage; while many states may have wished to project support for a sovereign state facing invasion, the system and identity narratives within a state may serve as a constraint on how far they can go. This creates the central tension that motivates our empirical analysis, and we develop our hypotheses around it in the following section.

A natural concern is that observed divergence between state media and official positions simply reflects internal incoherence rather than strategic communication. Indeed, weak democracies and authoritarian regimes should not be viewed as monolithic with respect to their internal political structure. Authoritarians face multiple factions they must balance---a general public and a cadre of elite insiders \citep{svolik_politics_2012}. These elite insiders are not just an arm of the ruler, but a set of actors with their own preferences and survival calculations that require continuous bargaining and cooptation on behalf of the autocrat \citep{bueno_de_mesquita_logic_2003}. Dictators must manage elite bargaining \citep{gandhi_political_2008}, and we extend this argument to posit that state-owned media, as an outlet of the regime, faces pressure from both the political leader and from elites involved with the running of the outlet itself, which may reflect different competing internal voices. Expressions of dissent even in the most constrained autocracies are consistent with different factional interests \citep{malesky_nodding_2010}. We argue that the signals sent by state-owned media should be considered important instruments of authoritarian communication, and in particular we should pay careful attention when those signals diverge from official positions. Whether state media diverges from formal positions because of top-down strategic choice or because of elite factional pressure, the divergence is analytically meaningful. While our data does not allow us to distinguish between these two mechanisms, our regression discontinuity design allows us to identify the causal effect of the invasion on state-media sentiment, independent of the diplomatic record. This means that divergence between the two is attributable to something other than common confounders, and we further develop our identification strategy in the methods section.



%%%%%%%%%%%%%%%%%%%%%%%%%%%%%% Original %%%%%%%%%%%%%%%%%%%%%%%%%%%%%%%%%%%%%%%%%%%%

Intuitively, the idea that state and independent media might react differently to a violation of an international border is not a novel concept. The two media sources might have very different mandates as well as funding sources. In addition, as a tool of the state, state-owned media is restricted by the states' decisions on how to make and signal alliance commitments. The international relations literature has grappled with several important questions related to alliances. Do states align with those most powerful and closest to them? Do states keep their friends closer, and their enemies closer? Upon deciding who to align with, how can small states, without powerful weapons and additional troops to send abroad, signal their commitment? The lack of concrete theory to help explain this variation in media response means that we must look to alliance, signaling, and media literature to help us begin to understand why we might see a difference in media sentiment. 

Weaker states often choose to balance in an effort to maintain their place in the system \citep{waltz_theory_2010}. When states do this, they are choosing the weaker side due to the perceived threat of the stronger state. States might also choose to bandwagon, that is join the stronger side \citep{waltz_theory_2010, walt_origins_2007}. States may choose to align with those who have the most aggregate offensive capabilities, who are located most closely, and who have the most aggressive perceived intentions. While bandwagoning might seem more attractive, Walt suggests that balancing is more common among states that matter and who have more capabilities \citep{walt_origins_2007}. 

While security concerns are likely to drive alliance preferences, states may also choose to join alliances that more closely adhere to their state views. Ideologically similar alliances are more common than alliances among countries of opposing ideological systems \citep{walt_origins_2007}. However, politicians might use ideological similarities to tout the strength of an alliance, though this level of closeness might be a facade, designed to inspire feelings of in-group favoritism. 

States may also enter into alliances designed to constrain. While some alliances are capability aggregation measures, meant to stand firm against external threats, alliances also may serve to tie squabbling partners together \citep{weitsman_dangerous_2004}. The alliance paradox suggests that allies are not only frequently unreliable, but also prone to confrontation. When states tether themselves by aligning with a powerful ally, they facilitate communication and mitigate potential hostilities, but this strategy constrains their partners. Tethering happens most often in relatively higher threat relationships; states that see themselves threatened by low levels of threat are more likely to hedge. States that are less threatened will "allow the hedging state to draw others into its sphere of influence as long as the course of action is not overly costly or too provocative to potential enemies" \citep{weitsman_dangerous_2004}. Small states located near powerful, threatening states might therefore choose to tether themselves to the regional pole, hoping to remain sovereign. 

Once states have decided on the best partners, how can they signal that they are committed? Signalling might involve words or deeds \citep{fuhrmann_signaling_2014}. Most political science literature focuses on how states might signal to rival states their level of capability and resolve \citep{fearon_signaling_1997, kertzer_resolve_2016}. Others have noted that signalling to allies is particularly difficult \citep{schelling_arms_1967}. Committed allies might make public statements or they might deploy troops or weapons to allied locations \citep{fuhrmann_signaling_2014, blankenship_trivial_2022}. 

By making public statements, countries could signal to their allies that they are committed to maintaining alliances. One of the most effective ways to broadcast these statements might be through the use of the media, including television, newspapers, and social media. This allows these statements to be public, helping to tie hands \citep{fearon_signaling_1997}. The news media is often used to set agendas and prime audiences as well as signal elite cues \citep{iyengar_news_2010, chong_framing_2007}. The news media can set agendas by identifying and publicizing issues for the mass public. The media can also prime by influencing the issues that the public pays attention to when they make judgements. Importantly, the media can help to transmit elite signals. There are two general types of media in most states: state-owned and private, or independent media. 

The distinction between state-owned and private media firms is unclear in the literature. Public choice theory suggests that government-owned media outlets “distort and manipulate information” to prop up the incumbent regimes, while privately owned and independent media sources offer alternative views with largely unbiased and accurate information \citep{djankov_who_2003}. State-owned media might also be more closely aligned with past authoritarian regimes \citep{moehler_whose_2011}. On the other hand, state-owned firms might provide more fair and balanced accounts, given their governmental affiliation and monetary backing \citep{djankov_who_2003}. Independent media, unhindered by government regulation may be more free to seek out the full story, or they may be backed by private firms that seek to provide a narrative more aligned with a desired end state. Privately owned media outlets might be more responsive to the public and open to opposing perspectives \citep{moehler_whose_2011}. Independent media in states that are more fully aligned with western values, with journalists interested in maintaining the freedom of the press that comes as a part of the international liberal order should be more likely to report more accurately. This will help them to maintain credibility as well as demonstrate their commitment to showing what is going on in the world. 

Within the sphere of state-owned media, there are two primary strategic threads discussed in the literature. The dominant strand has focused on the use of state-owned media in communication with domestic audiences. For example, \citet{mccombs_agenda-setting_2002} emphasizes the power of the news media to focus public attention and shape perspectives. Careful framing of issues within media can have a dramatic impact on public support for government actions \citep{gershkoff_shaping_2005}. Political parties wield significant control over their followers' opinions and can wield this influence to effect dramatic changes in support for policies \citep{slothuus_how_2021}. This is in conjunction with an emphasis within the broader discipline on domestic affairs \citep{tomz_public_2020, de_mesquita_testing_2004}. However, there has been a recent revitalization of research into the strategic use of media signalling to influence international audiences. 

Recent work has highlighted the extent to which savvy autocracies maintain power through the strategic manipulation of information--- management of state-owned media has become sophisticated, especially during crises \citep{guriev_spin_2022}. At the same time, there are important differences both across and within autocracies. Across autocracies, the underlying resource wealth of a state has been shown to have a profound impact on levels of media freedom \citep{egorov_why_2009}, which might help explain different incentives for state media manipulation and points to a need to control for the underlying media environment within a state. In particular, \citet{weeks_autocratic_2008} argues that different autocracies have systematic differences in their behavior in the international sphere. We argue that a result of these systematic differences is variation in their capacity for credible formal signaling. Those states with less capability for formal signaling may have an incentive to supplement their formal signals through informal channels that can convey preferences that formal diplomatic constraints make difficult to express directly. We argue that state-owned media is a plausible channel for this communication for two reasons. First, there is historical reliance by many countries in our sample on state-owned media for international signaling \citep{cohen_theatre_1987}. Second, the explicit rise in efforts by western countries to systematically study state media \citep{manor_digitalization_2019, usic_ic_2024} gives senders reason to believe that the signal will be received and correctly interpreted.

A natural concern is that observed divergence between state media and official positions simply reflects internal incoherence rather than strategic communication. Indeed, weak democracies and authoritarian regimes should not be viewed as monolithic with respect to their internal political structure. Authoritarians face multiple factions they must balance, a general public and a cadre of elite insiders \citep{svolik_politics_2012}. These elite insiders are not just an arm of the ruler, but a set of actors with their own preferences and survival calculations that require continuous bargaining and cooptation on behalf of the autocrat \citep{bueno_de_mesquita_logic_2003}. Dictators must manage elite bargaining \citep{gandhi_political_2008}, and we extend this argument to posit that state-owned media, as an outlet of the regime, faces pressure from both the political leader, but also from elites involved with the running of the outlet itself which may reflect different competing internal voices. Expressions of dissent even in the most constrained autocracies are consistent with different factional interests \citep{malesky_nodding_2010}. We argue that the signals sent by state-owned media should be considered important instruments of authoritarian communication and in particular, we should pay careful attention when the signals of state-owned media diverge from official positions. Whether state media diverges from formal positions because of top-down strategic choice or because of elite factional pressure, the divergence is analytically meaningful. While our data does not allow us to distinguish between these two mechanisms, our regression discontinuity design allows us to identify the causal effect of the invasion on state-media sentiment, independent of the diplomatic record. This means that divergence between the two is attributable to something other than common confounders, and we further develop our identification strategy in our methods section.

%While there is significant variation across the type of government in our dataset, we argue that in this particular crisis, the best predictor of their behavior is their alignment within the international sphere.


Governments have a strong incentive to capture media. \citet{gehlbach_government_2014} provides a workhorse theoretical model for strategic state-owned media. They argue that state-owned media is a strategic tool used by government's to advance a specific agenda and as a result, bias in state-owned media is generally higher than in independent media. While they argue that, with sufficient incentive to advance their agenda, the government is more likely to engage in capture of independent media, either directly or indirectly through paying off media, in a crisis like the invasion of Ukraine, either strategy is unlikely to be immediately actionable in the short-term. Building upon this logic, we argue that we should expect to see a strong divergence between state-owned media coverage and independent media coverage initially around the invasion as state-owned media responds strategically and before governments can begin to influence independent media coverage. 

Importantly, state-owned media is effective in shaping audience opinions in weak democracies. \citet{enikolopov_media_2011} provides empirical evidence of the impact of independent media coverage. They show that exposure to independent media dramatically reduces vote-share for the ruling party in the 1999 Russian election. This highlights the extent to which state-owned media framing can matter in the absence of exposure to independent media--- it serves not just as a tool for mobilization of existing preferences, but can serve to shape and change attitudes. Importantly, they argue that media plays a powerful role in persuasion in countries with weak democratic institutions, which characterizes most of the countries in our sample. 

There is evidence that state-owned media has been used as a geopolitical instrument in the region. \citet{peisakhin_electoral_2018} use variation in exposure to Russian state TV to examine support for the Pro-Russian candidates in Ukraine's 2014 elections. They find that exposure increases support for Pro-Russian candidates, but primarily amongst citizens already inclined to support those candidates. This has two interesting implications for our paper; first, it provides additional evidence of the importance of the strategic use of state-owned media in the region, and second, it highlights that state-owned media can serve as an important geopolitical instrument.

Within the region, \citet{oates_television_2006} argues that state-owned media serves most directly as a tool for elites to engage in agenda setting. She highlights how effective it has been in managing political discourse and short-circuiting efforts to promote democratic institutions. This suggests that there are benefits to state-owned media strategically manipulating sentiment around an event like the invasion of Ukraine. This is particularly important within the context of the contemporary media ecosystem with readily available online spheres. Careful use of disinformation can provide advantages to repressive regimes seeking to manipulate audience opinions \citep{oates_russian_2016}.

The relationship between media and foreign policy has been examined through several lenses. The foundational text on media and foreign policy, \citet{robinson_cnn_2002} explored the idea of, and ultimately found mixed evidence for, a "CNN effect", whereby media coverage of crises forces government to intervene. Instead, he found that media influence is only likely when policy is uncertain and coverage is both critical of government and empathetic of those suffering during the crisis. Once policy is settled, he argues that media largely follows elite cues. We reframe this; we argue that rather than simply following elite cues or driving foreign policy, states with access to state-owned media strategically weaponize media as foreign policy.

\citet{entman_projections_2004} introduced a cascade model of media, with frames originating with the executive and filtering down through elites, media, and finally to the public, with friction at each level within a democracy. We build on his framework of media framing as a modern foreign policy instrument by extending his work to autocracies and weak democracies, where we argue that there is less friction at each level, especially with respect to state-owned media. This is driven by a combination of editorial processes, which explicitly work to build on official state positions, and journalistic incentives, where journalists face potentially harsh sanctioning for challenging the state position.

\citet{cohen_theatre_1987} emphasizes the role of diplomatic signaling through the media. He argues that these media signals are are ambiguous by design, allowing deniability, but are well understood by the intended recipients (other states)---- a dynamic that has only intensified as foreign governments now systematically monitor state media output \citep{manor_digitalization_2019, bjola_revisiting_2018, usic_ic_2024}. Indeed, during the Cold War, western analysts meticulously studied the state-owned media of the USSR and its satellites to understand the Soviet foreign policy. Soviet officials strategically used this fact to signal to the West in ways that they could not do through formal channels \citep{cohen_theatre_1987}. This work serves as a precursor to \citet{miskimmon_strategic_2013}, which introduces the concept of strategic narratives, which states leverage through media as a core instrument of foreign policy and serve to shape how states project their narrative to the world. Their framework establishes how actors form and project narratives at three levels; system, identity and issue. 

The need for informal signaling channels is particularly acute for autocracies. Formal signaling theory establishes that credible diplomatic signals require that the signaling state incur costs that make bluffing irrational, either by tying hands through audience costs or by sinking material costs \citep{fearon_signaling_1997}. Democratic governments can generate audience costs through open political competition, where visible domestic dissent or support provides external observers with reliable information about a government's true resolve \citep{schultz_democracy_2001}. Autocracies and weak democracies, which suppress dissent and coerce apparent consensus, face structural credibility deficits in formal channels: their official diplomatic statements are less informative precisely because external observers know they carry little or no domestic political cost. State media, in this context, becomes a more credible--- albeit still deniable--- channel because foreign governments monitor it as a less-controlled reflection of elite preferences \citep{bjola_revisiting_2018}. The digitalization of public diplomacy has formalized this dynamic. State media is now systematically tracked by foreign intelligence and diplomatic services as informative about the actual policy preferences of the sending government \citep{manor_digitalization_2019, usic_ic_2024}.

We extend the framework proposed by \citet{miskimmon_strategic_2013} by considering the constrained nature of international signaling due to a combination of trade and security integrations. We posit that strategic narrative projection operates at both formal levels (UN votes, official diplomatic statements) and more informal levels (state-owned media). We argue that the invasion of Ukraine triggers crises at all three levels of the framework in \citet{miskimmon_strategic_2013}. While the invasion serves foremost as an issue narrative contest, the implications of a major land war in Europe, an international order shaking crisis, triggers both system and identity pressures for post-Soviet states. The invasion is a direct challenge to the established post-Cold War system, challenging norms around the use of force and wars of territorial expansion. At the identity level, it challenged state positions--- its values, its alliances and its place in the world. For post-Soviet states that found themselves caught between Russian security agreements and Western integration, this introduced enormous identity issues. That this crisis provoked challenges around all three levels of strategic narratives introduced a tension that states needed to manage; while many states may have wished to project support for a sovereign state facing invasion, the system and identity narratives within a state may serve as a constraint to how far it can go. This serves as the crux of our third hypothesis--- states with Russian security agreements were forced to rely on issue-level media signals to express preferences that identity- and system-level constraints made impossible to communicate through formal diplomatic channels.


\newpage

\printbibliography

%\bibliographystyle{apalike}
%\bibliography{references.bib}

\end{document}